% Generated into workshop-exercises.tex by build.py. Edit this file, not the output.
\documentclass[9pt]{extarticle}
\usepackage[a4paper,top=10mm,bottom=13mm,left=11mm,right=11mm,footskip=6mm]{geometry}
\usepackage[T1]{fontenc}
\usepackage[utf8]{inputenc}
\usepackage[default]{lato}
\usepackage[scaled=0.9]{inconsolata}
\usepackage{textcomp}
\usepackage{microtype}
\DisableLigatures{encoding=*, family=*}
\usepackage{xcolor,graphicx,multicol,qrcode,fancyhdr,lastpage}
\usepackage[most]{tcolorbox}
\usepackage[hidelinks]{hyperref}
\hypersetup{pdftitle={AI in the Workplace: Copilot Exercises},pdfauthor={Dr Matt Mort and Mr Matt Hayden, Cardiff University}}

% Copy and paste: real spaces, Unicode mapping, no hyphenation.
\input glyphtounicode
\pdfgentounicode=1
\pdfinterwordspaceon
\hyphenpenalty=10000 \exhyphenpenalty=10000
\raggedright
\setlength{\parindent}{0pt}
\setlength{\columnsep}{6mm}

\definecolor{ink}{HTML}{1D1D1F}
\definecolor{muted}{HTML}{4A4F55}
\definecolor{accent}{HTML}{0071E3}
\definecolor{tint}{HTML}{F2F7FD}
\definecolor{rule}{HTML}{C9DDF5}
\definecolor{gold}{HTML}{FF8A1F}
\color{ink}

\pagestyle{fancy}\fancyhf{}
\renewcommand{\headrulewidth}{0pt}
\fancyfoot[L]{\footnotesize\color{muted}Dr Matt Mort and Mr Matt Hayden, Cardiff University \quad\textbullet\quad \href{@@URL@@}{@@URL@@}}
\fancyfoot[R]{\footnotesize\color{muted}Page \thepage\ of \pageref{LastPage}}

% Prompt box: everything inside is the prompt and nothing else.
\newtcolorbox{pbox}{enhanced,colback=tint,colframe=rule,boxrule=0.4pt,arc=1.5mm,
  left=2mm,right=2mm,top=1mm,bottom=1mm,before skip=0.8mm,after skip=1.2mm,fontupper=\pboxfont\raggedright}
\newcommand{\pboxfont}{\normalsize}
\newcommand{\gap}{\par\vspace{0.5ex}}
\newcommand{\bul}[1]{\par\hangindent=1em\hangafter=1\textbullet\ #1\par}
\newcommand{\codeline}[1]{\par{\ttfamily\footnotesize #1}\par}

\newcommand{\exskip}{4.5mm}
\newcommand{\ex}[2]{\par\vspace{\exskip}\needspace{4\baselineskip}%
  {\setlength{\fboxsep}{2pt}\colorbox{accent}{\makebox[\widthof{\bfseries\small M365}][c]{\color{white}\bfseries\small\strut #1}}\hspace{1.5mm}{\bfseries\large #2}}\hspace{2.5mm}}
\newcommand{\note}[1]{{\small\color{muted}#1}\par}
\newcommand{\sub}[1]{\par\vspace{1.5mm}{\small\bfseries #1}\par}
\usepackage{needspace,calc}

\newcommand{\band}[2]{%
  \begin{tcolorbox}[enhanced,colback=ink,colframe=ink,arc=2mm,left=4mm,right=3mm,top=2.5mm,bottom=2.5mm]
  \begin{minipage}[c]{0.8\linewidth}\color{white}\raggedright
    {\fontsize{17}{20}\selectfont\bfseries #1}\par\vspace{1.2mm}
    {\small #2}
  \end{minipage}\hfill
  \begin{minipage}[c]{0.15\linewidth}\raggedleft
    \colorbox{white}{\qrcode[height=1.75cm]{@@URL@@}}
  \end{minipage}
  \end{tcolorbox}\vspace{-1mm}}

\begin{document}

\band{AI in the Workplace: Copilot Exercises}{%
  Open Copilot at \href{https://m365copilot.com}{\textbf{https://m365copilot.com}}, start a new chat, copy a prompt from a blue box, paste it in and send.\par\vspace{0.6mm}
  Full workshop, any time, at home or at work: \href{@@URL@@}{\textcolor{gold}{\textbf{@@URL@@}}}}

\ex{01}{Explain it like I am 5}
\note{Replace [TOPIC], or upload a document, such as a policy, and ask Copilot to explain it.}
@@progressive-learning@@

\ex{02}{Roleplay a difficult conversation}
\note{Copilot stays in character. Reply as yourself.}
@@roleplay-conversation@@

\ex{03}{Devil's advocate}
\note{Replace [SCENARIO HERE] with your own situation, or upload a document instead.}
@@devils-advocate@@

\ex{04}{Document to infographic}
\note{Click \textbf{+}, \textbf{Upload document}, attach any document, then paste. No file? Use a link or topic in FOCUS ON.}
@@infographic-prompt@@

\ex{05}{Professional persona}
\note{Try both. Same question, new chat, different right answers.}
@@role-prompt-medical@@
@@role-prompt-comedian@@

\ex{06}{Format and structure}
\note{One prompt, three formats in a single reply.}
@@format-prompt-merged@@

\ex{07}{Context: know your audience}
\note{Then ask Copilot to explain its choices, and give it feedback.}
@@context-prompt@@

\newpage
\renewcommand{\exskip}{1.5mm}

\band{Memory and Bonus Prompts}{%
  Full instructions and examples: \href{@@URL@@}{\textcolor{gold}{\textbf{@@URL@@}}}}

\ex{08}{Copilot memory}
\note{Rules saved to memory apply to every future chat. To view them: three dots (top right), \textbf{Settings}, \textbf{Personalisation}.}
@@memory-regional@@
@@memory-honesty@@

\ex{09}{Temporary chat}
\note{Click the temporary chat icon \raisebox{-0.35\height}{\includegraphics[height=5mm]{../copilot-temporary-chat-icon.png}} (top right). Not saved and creates no memories, but it is not private.}
\vspace{1.2mm}

\ex{B1}{Copy the style of a document you like}
\note{Attach a document or image whose style you want to reuse. PDFs work best.}
@@clone-style-prompt@@

\ex{B2}{Turn a good chat into a reusable prompt}
\note{Give Copilot feedback until the output is right, then paste this.}
@@feedback-better-prompt@@

\ex{B3}{Handover: long chats make more mistakes}
\note{Paste this into the long chat. Copy the reply, start a new chat and paste it in.}
@@context-rot-handover@@

\ex{B4}{Treat your chat history as a resource}
\sub{Build a full profile of yourself}
@@chat-history-profile@@
\sub{Find prompting habits you repeat}
@@chat-history-habits@@
\sub{Spot tasks you keep redoing by hand}
@@chat-history-tasks@@
\sub{Capture your writing style}
@@chat-history-style@@

\ex{M365}{Devil's advocate on your own work}
\note{Microsoft 365 Copilot licence only. It reads your emails and meetings.}
@@devils-advocate-work@@

\end{document}
